
在本章中，我们将要定义\trActionFunctional{}的诸\trCriticalPoint{}的\trIndex。

按照传统，\trMaslovIndex{}是一个与\trLoop{}相伴的整数，这里的\trLoop{}取在 $\RR^{2n}$ 的全部\trLagrange{}子空间所成的\trGrassmann\trManifold{} $\Lambda_n$ 之中。这是因为这个\trGrassmann\trManifold{}的\trFundamentalGroup{}与 $\ZZ$ \trIsomorphic（例如参见\cref{ex:ch14-040-lagrangian-grassmannian}，或书 [50]），而归根结底，则是因为酉群 $\mathrm{U}(n)$ 的\trFundamentalGroup{}同样与 $\ZZ$ \trIsomorphic。我们在此要谈论的这个\trIndex，则是一个与\trSymplecticGroup{} $\mathrm{Sp}(2n)$ 中的道路相伴的整数。后者可收缩到 $\mathrm{U}(n)$ 上，从而也具有一个无限循环的\trFundamentalGroup。在这个框架之下，我们也把它称作\trConleyZehnderIndex（欲知文献，参见 [67, 66]、[49]——我们曾由之汲取灵感——以及近期的论文 [43]）。

\section{迈向\trIndex{}的定义}
\label{sec:toward-definition-index}

本节的框架仍是前一章的框架；我们有一个时间依赖的\trHamiltonian{}量
\[
H\colon W\times S^1\longrightarrow\RR
\]
以及方程 $\dot{x}=X_t(x)$ 的一个\trPeriodic{}为 $1$ 的解 $x(t)$，我们假设它是一条\trContractible{}的\trLoop。我们已经看到，这些\trLoop{}正是 $\scr{A}_H$ 的诸\trCriticalPoint。于是，我们想为这些\trCriticalPoint{}定义一个\trIndex。

我们将分三步来做这件事。我们从一个\trNondegenerate{}（且\trContractible{}）的\trTrajectory{} $x$ 出发。

\begin{enumerate}[label=(\arabic*)]
\item 我们赋予它一条\emph{道路} $t\mapsto A(t)$，由\emph{\trSymplectic}矩阵构成，使得 $A(0)=\mathrm{Id}$ 且 $1\notin\operatorname{Spec}A(1)$。
\item 对每一条这样的道路，我们都赋予一条道路 $\gamma\colon[0,1]\to S^1$，使得 $\gamma(0)=1$ 且 $\gamma(1)=\pm1$。
\item 再对每一条这样的道路（其取值在 $S^1$ 中），我们赋予一个整数。
\end{enumerate}

给定一个\trNondegenerate{}的\trCriticalPoint{} $x$，我们便由此赋予了它一个整数 $\mu(x)\in\ZZ$。

\subsection{第一步}
\label{subsec:first-step}

我们已经固定了\trPeriodic{}\trTrajectory{} $x(t)=\varphi^t(x(0))$。我们选取 $T_{x(0)}W$ 的一组\trSymplectic{}基
\[
Z(0)=\bigl(Z_1(0),\,\dotsc,\,Z_{2n}(0)\bigr)\text{。}
\]
回想一下（参见\cref{sec:hamiltonian-vector-fields}），$\varphi^t$ 保持 $\omega$。特别地，线性映射 $T_{x(0)}\varphi^1$ 在基 $Z(0)$ 下的矩阵是一个\trSymplectic{}矩阵 $A(1)\in\mathrm{Sp}(2n)$（关于\trSymplecticGroup{} $\mathrm{Sp}(2n)$，参见\cref{sec:symplectic-group}），而 $1$ 不是 $A(1)$ 的\trEigenvalue，因为解 $x$ 已被假定为\trNondegenerate{}的。

由下面这条经典定理（例如参见 [75]），存在 $T_{x(t)}W$ 的一族\trSymplectic{}基
\[
Z(t)=\bigl(Z_1(t),\,\dotsc,\,Z_{2n}(t)\bigr)\text{，}
\]
它关于 $t$ 是 $\cC^\infty$ 的。

\begin{theorem}\label{thm:disk-bundle-trivialization}
对每个连续映射
\[
\psi\colon D^k\longrightarrow W\text{，}
\]
\trSymplectic\trFiberBundle{} $\psi^\star TW$ 都存在\trTrivialization，并且其一切\trTrivialization{}都相互\trHomotopic。\noqed
\end{theorem}

我们将不证明这条定理；它由\trDisk{} $D^k$ 的\trContractible{}性得出。例如，一切\trTrivialization{}的\trHomotopic{}等价性，可由“两个\trTrivialization{}被一个映射 $D^k\to\mathrm{Sp}(2n)$ 相联系”这一事实得出，而且它只用到 $\mathrm{Sp}(2n)$ 是道路\trConnected{}的这一事实（\cref{thm:sp-polar-retraction}）。

言归正传。我们的\trPeriodic{}\trTrajectory{} $x$ 是\trContractible{}的，我们把它扩张成一个映射 $u\colon D^2\to W$；由这条定理，我们导出一个\trTrivialization，也就是说，一个在 $z\in D^2$ 上定义的\textbf{\trSymplecticFrame{}}（symplectic frame）$Z(z)$，特别是沿 $x$ 的一个\trSymplecticFrame。

这样一来，对每个 $t$，我们都可以考虑线性映射 $T_{x(0)}\varphi^t$ 在基 $Z(0)$（$T_{x(0)}W$ 的）与 $Z(t)$（$T_{x(t)}W$ 的）之下的矩阵 $A(t)$，而映射
\[
t\longmapsto A(t)
\]
是一条道路，其中 $A(0)=\mathrm{Id}$，并且 $A(1)$ 不具有\trEigenvalue{} $1$。

于是让我们设
\[
\mathrm{Sp}(2n)^\star=\{A\in\mathrm{Sp}(2n)\mid\det(A-\mathrm{Id})\neq 0\}
\]
并设 $\scr{S}$ 为把 $\mathrm{Id}$ 连到 $\mathrm{Sp}(2n)^\star$ 中一个元素的那些道路所成的空间：
\[
\scr{S}=\bigl\{\gamma\colon[0,1]\longrightarrow\mathrm{Sp}(2n)\mid\gamma(0)=\mathrm{Id}\ \text{且}\ \gamma(1)\in\mathrm{Sp}(2n)^\star\bigr\}\text{。}
\]

若 $Z'$ 是沿 $x$ 的另一个\trSymplecticFrame，则道路 $A(t)$ 与 $A'(t)$ 在 $\scr{S}$ 中\trHomotopic。\footnote{译注：原书此处称，凡沿 $x$ 的另一\trSymplecticFrame{}皆给出与 $A(t)$ \trHomotopic{}的道路；然沿\trLoop{}任意更换\trSymplecticFrame{}并不能保证这一点——沿\trLoop{}绕行一周的\trSymplecticFrame{}变更（给道路左乘一枚转完 $2\pi$ 的旋转回路，如 $A'(t)=R(-2\pi t)A(t)$，其中 $R(\theta)$ 表转角为 $\theta$ 的旋转），所得道路便与 $A(t)$ 在 $\scr{S}$ 中不\trHomotopic，且二者相应的\trIndex{}恰差 $2$。故此句应补上条件：$Z'$ 可延拓至所选的\trCappingDisk{} $u$（或与原\trSymplecticFrame{} $Z$ \trHomotopy）。此处照原书排印，存疑备案。——译者注}更一般地，若 $v$ 是 $x$ 到\trDisk{} $D^2$ 的另一个扩张，并且 $Z'$ 是沿 $v$ 的一个\trTrivialization，那么通过把 $u$ 与 $v$ \trGluing{}起来，我们便能构造一个映射 $w\colon S^2\to W$（如在\cref{subsec:action-functional-b}中；参见\cref{fig:disk-gluing-sphere}）。于是我们有：
\begin{itemize}
\item 由\trTrivialization{} $Z$ 沿 $u$ 给出的道路 $A(t)$；
\item 由 $Z'$ 沿 $v$ 给出的道路 $A'(t)$；
\item 还有一条道路 $A''(t)$，由一个\trTrivialization{} $Z''$ 沿 $w$ 给出——多亏了\cref{thm:symplectic-trivialization}，
\end{itemize}
而这三条道路在 $\scr{S}$ 中\trHomotopic。

我们由此完成了第一步，也就是说，赋予了\trTrajectory{} $x$ 一条 $\scr{S}$ 中的道路，它在\trHomotopy{}意义下是唯一的。

\begin{remark}\label{thm:autonomous-hessian-path}
对于\trAutonomous{}\trHamiltonian{}量 $H$ 的一个\trCriticalPoint{} $x$ 的情形，我们已经构造过道路 $t\mapsto A(t)$（在\cref{sec:hamiltonian-vector-fields}中，以及\cref{ex:ch14-008-jacobian-flow}里）：它就是
\[
A(t)=e^{tJ_0\,\operatorname{Hess}_x}
\]
其中 $\operatorname{Hess}_x$ 是 $H$ 在 $x$ 处的\trHessian，为简单起见，我们将以 $S$ 记之。特别地，若 $S$ \trNonsingular{}的且 $\|S\|<2\pi$，则 $J_0S$ 不具有\trEigenvalue{} $2ik\pi$，从而 $e^{J_0S}$ 不具有\trEigenvalue{} $1$。
\end{remark}

\subsection{第二步（开端）}
\label{subsec:second-step-beginning}

至于第二步，对一个道路 $A(t)\in\scr{S}$，我们必须赋予一条道路
\[
\gamma\colon[0,1]\longrightarrow S^1\qquad\text{其中}\ \gamma(0)=1\ \text{且}\ \gamma(1)=\pm1\text{。}
\]
为了做到这一点，我们需要\trSymplecticGroup{}的若干性质，也就是第 5 章中见过的那些，外加一些补充细节。

在\cref{subsec:polar-decomposition}中我们曾说过，存在一个映射 $\mathrm{Sp}(2n)\to S^1$，它在\trFundamentalGroup{}上诱导一个\trIsomorphism。在本节中，我们将给出关于这样一个映射 $\rho$ 之性质的更多细节。此处给出的构造，基于 \trSalamon{}与\trZehnder{}在 [67] 中对 [38] 的运用。我们将给出一个略有不同（尽管等价）的描述，并且把我们所用到的那些线性代数结果的详细证明写下来。

\begin{theorem}[{[67]}]\label{thm:rho-map-properties}
对每个 $n\in\NN^\star$，都存在一个连续映射
\[
\rho\colon\mathrm{Sp}(2n)\longrightarrow S^1
\]
满足下列诸性质：
\begin{enumerate}[label=(\arabic*)]
\item 自然性：若 $A$ 与 $T\in\mathrm{Sp}(2n;\RR)$，则
\[
\rho(TAT^{-1})=\rho(A)\text{。}
\]
\item 乘积：若 $A\in\mathrm{Sp}(2m;\RR)$ 与 $B\in\mathrm{Sp}(2n;\RR)$，则
\[
\rho\begin{pmatrix}A&0\\0&B\end{pmatrix}=\rho(A)\rho(B)\text{。}
\]
\item 行列式：若 $A\in\mathrm{U}(n)=\mathrm{Sp}(2n;\RR)\cap\mathrm{O}(2n)$，则
\[
\rho(A)=\det\nolimits_{\CC}(X+iY)\text{，}\qquad\text{其中 }A=\begin{pmatrix}X&-Y\\Y&X\end{pmatrix}\text{。}
\]
特别地，$\rho$ 诱导一个\trIsomorphism
\[
\rho_\star\colon\pi_1(\mathrm{Sp}(2n);\RR)\longrightarrow\pi_1(S^1)=\ZZ\text{。}
\]
\item 规范化：若 $\operatorname{Spec}(A)\subset\RR$，则
\[
\rho(A)=(-1)^{m_0/2}\text{，}
\]
其中 $m_0$ 是诸负实\trEigenvalue{}的总\trMultiplicity。
\item $\rho({}^tA)=\rho(A^{-1})=\overline{\rho(A)}$。
\end{enumerate}
\end{theorem}

构造这样一个映射 $\rho$ 是相当技术性的。我们将在本章的附录（§7.3）中给出细节。尽管如此，不妨先说一句：对一个诸\trEigenvalue{}两两相异的矩阵 $A$ 而言，这“不过”就是
\[
\rho(A)=(-1)^{m_0/2}\prod_{\substack{\lambda\in\operatorname{Spec}(A)\cap S^1\\ \operatorname{Im}(\lambda)>0}}\lambda^{\operatorname{sign}\operatorname{Im}\,\omega(\overline{X},X)}
\]
（对每一个 $\lambda$，我们都已选定了一个特征向量 $X$）。

\subsection{子集 $\mathrm{Sp}(2n)^\star$}
\label{subsec:sp-star-subset}

让我们回到对\trPeriodic{}\trTrajectory{} $x$ 的\trIndex{} $\mu$ 的构造上来。

\textbf{不具有\trEigenvalue{} $1$ 的诸变换。}

我们需要考虑的那些\trSymplectic{}变换不具有\trEigenvalue{} $1$（这就是我们的\trNondegenerate{}性假设）。这就是为什么我们考虑\trOpen{}集
\[
\mathrm{Sp}(2n)^\star=\{A\in\mathrm{Sp}(2n)\mid\det(A-\mathrm{Id})\neq 0\}\text{，}
\]
它是由
\[
\Sigma=\{A\in\mathrm{Sp}(2n)\mid\det(A-\mathrm{Id})=0\}
\]
所定义的\trHypersurface{} $\Sigma$ 的\textbf{补集}。

这就使我们向诸\trLagrange{}的\trGrassmann\trManifold{}中\trMaslovIndex{}的通常定义靠近了一步。

显而易见，$\mathrm{Sp}(2n)^\star$ 不可能是\trConnected{}的，因为它是由 $\det(A-\mathrm{Id})>0$ 与 $\det(A-\mathrm{Id})<0$ 所定义的两个不交\trOpen{}集的并。还要注意，$\det(A-\mathrm{Id})$ 的符号与乘积 $\prod(\lambda_i-1)$ 的符号相同，其中 $\lambda_i$ 取遍 $A$ 的那些正实\trEigenvalue。

\textbf{例：$\mathrm{Sp}(2)$。}

在 $n=1$ 的情形，$\mathrm{Sp}(2)=\mathrm{SL}(2;\RR)$。这个群是由满足 $ad-bc=1$ 的实矩阵 $\begin{pmatrix}a&b\\c&d\end{pmatrix}$ 所组成的（实仿射的）\trQuadric，一个维数为 $3$ 的\trSubmanifold。\trPolarDecomposition{}使它收缩到由矩阵 $\begin{pmatrix}\cos\theta&-\sin\theta\\ \sin\theta&\cos\theta\end{pmatrix}$ 所组成的\trCircle{}上（这就是群 $\mathrm{U}(1)$）。由具有\trEigenvalue{} $1$ 的矩阵组成的这个\trHypersurface{} $\Sigma$，则由满足 $ad-bc=1$ 且 $a+d=2$ 的诸 $\begin{pmatrix}a&b\\c&d\end{pmatrix}$ 构成。它在点 $\mathrm{Id}$ 之外是\trSmooth{}的。$\Sigma$ 的\textbf{补集}分别由
\[
\operatorname{tr}A>2\qquad\text{以及}\qquad \operatorname{tr}A<2
\]
所定义的两个\trOpen{}集组成。人们易于看出这两个\trOpen{}集都是\trConnected{}的，因为：
\begin{itemize}
\item \textbf{\trTrace{}}（trace）大于 $2$ 的那些矩阵，正是与
\[
\begin{pmatrix}\lambda&0\\0&\lambda^{-1}\end{pmatrix}\quad(\text{其中 }\lambda>0)
\]
相似的矩阵。\footnote{译注：按“当且仅当”的读法，$\lambda>0$ 须排除 $\lambda=1$（$\lambda=1$ 时该矩阵就是 $\mathrm{Id}$，迹为 $2$，本不属于 $\operatorname{tr}A>2$ 一侧）；同理，下一款的旋转矩阵须排除 $\theta\in2\pi\ZZ$（此时旋转也是 $\mathrm{Id}$）；此处照原书排印。——译者注}
\item \trTrace{}小于 $2$ 的那些矩阵，正是与
\[
\begin{pmatrix}\cos\theta&-\sin\theta\\ \sin\theta&\cos\theta\end{pmatrix}\text{，}\quad \begin{pmatrix}\lambda&0\\0&\lambda^{-1}\end{pmatrix}\quad(\text{其中 }\lambda<0)\text{，}\quad\text{或}\quad \begin{pmatrix}-1&a\\0&-1\end{pmatrix}\cdots
\]
相似的矩阵——全都是能够毫无问题地连到 $-\mathrm{Id}$ 的矩阵。
\end{itemize}

\noindent\textbf{注记。}映射 $A\mapsto\rho(A)$ 并不等于 $U$ 的复行列式，其中 $A=US$ 是\trPolarDecomposition（参见\cref{subsec:polar-decomposition}）。这里有一个反例：取
\[
U=\begin{pmatrix}\cos\pi/3&-\sin\pi/3\\ \sin\pi/3&\cos\pi/3\end{pmatrix}\qquad\text{以及}\qquad S=\begin{pmatrix}4&0\\0&1/4\end{pmatrix}\text{。}
\]
于是 $\operatorname{tr}(A)>2$，从而根据\cref{thm:rho-map-properties}的性质 (4)，$\rho(A)=1$。

\textbf{$\mathrm{Sp}(2n)^\star$ 的\trTopology}

下述命题归功于\trConley{}与\trZehnder。

\begin{proposition}[{[21]}]\label{thm:sp-star-components}
\trOpen{}集 $\mathrm{Sp}(2n)^\star$ 具有两个\trConnectedComponent。它们各自到 $\mathrm{Sp}(2n)$ 中的包含，都在\trFundamentalGroup{}上诱导零\trHomomorphism。
\end{proposition}

\begin{proof}
我们已经说过，$\mathrm{Sp}(2n)^\star$ 不是\trConnected{}的。我们尚须证明，两个子空间
\[
\mathrm{Sp}(2n)^{+}=\{A\in\mathrm{Sp}(2n;\RR)\mid\det(A-\mathrm{Id})>0\}
\]
与
\[
\mathrm{Sp}(2n)^{-}=\{A\in\mathrm{Sp}(2n;\RR)\mid\det(A-\mathrm{Id})<0\}
\]
都是道路\trConnected{}的。这一事实的证明基于下面这条引理。

\begin{lemma}\label{thm:distinct-positive-eigenvalues}
设 $A\in\mathrm{Sp}(2n)^\star$。则存在 $\mathrm{Sp}(2n)^\star$ 中的一条道路，把 $A$ 连到一个矩阵 $B$，使得 $B$ 的诸\trEigenvalue{}两两相异，并且 $B$ 恰有零个（若 $A\in\mathrm{Sp}(2n)^{+}$）或两个（若 $A\in\mathrm{Sp}(2n)^{-}$）正实\trEigenvalue。
\end{lemma}

我们先承认这条引理（暂时地），用它来证明：$\mathrm{Sp}(2n)^\star$ 中的每个矩阵 $A$ 都能经由 $\mathrm{Sp}(2n)^\star$ 中的一条道路，连到矩阵
\[
\begin{cases}
-\mathrm{Id}\in\mathrm{Sp}(2n)^{+} & \text{若 }A\in\mathrm{Sp}(2n)^{+}\text{，}\\[8pt]
\left(\begin{array}{cc|c}
2 & 0 & \\
0 & 1/2 & 0\\
\hline
\multicolumn{2}{c|}{0} & -\mathrm{Id}
\end{array}\right) & \text{否则。}
\end{cases}
\]
我们先把 $A$ 连到引理所容许的一个矩阵 $B$；接下来，我们还须把 $B$ 连到上面两个矩阵之一。我们选取 $B$ 的一组由特征向量构成的基，它既\trSymplectic{}又在复共轭下不变。\footnote{译注：此语宜理解为：选取\trAdapted{}于复共轭与辛配对的（复）谱基——$B$ 是实矩阵，对非实的\trEigenvalue{} $\lambda$ 取 $X_{\overline{\lambda}}=\overline{X_{\lambda}}$，并按\trSymplecticForm{}配对归置诸特征向量，“既\trSymplectic{}又在复共轭下不变”即就此而言；不宜读作由实特征向量组成的基（非实的\trEigenvalue{}没有实特征向量）。此处照原书排印。——译者注}由于 $B$ 的诸\trEigenvalue{}两两相异，我们可以把与 $\lambda$ 相伴的那个特征向量称作“这个”特征向量 $X_\lambda$。对 $B$ 的每一个不是正实数的\trEigenvalue{} $\lambda$，我们选取一条道路 $\lambda(s)$，它把 $\lambda$ 连到 $-1$ 而避开 $1$，并且使得：若 $\lambda(s)$ 是为 $\lambda$ 选定的道路，那么：
\begin{itemize}
\item 为 $\overline{\lambda}$ 选定的道路是 $\overline{\lambda}(s)\ldots$
\item 而为 $1/\lambda$ 选定的道路是 $1/\lambda(s)$。
\end{itemize}
我们用下式来定义道路 $s\mapsto B(s)$：
\[
B(s)\cdot X_\lambda=
\begin{cases}
\lambda(s)X_\lambda & \text{若 }\lambda\notin\RR_+\text{，}\\
\lambda X_\lambda & \text{若 }\lambda\in\RR_+\text{。}
\end{cases}
\]
所得的矩阵 $B(s)$ 是实的且属于 $\mathrm{Sp}(2n)^\star$；$B(0)=B$，并且
\[
B(1)\cdot X_\lambda=
\begin{cases}
-X_\lambda & \text{若 }\lambda\notin\RR_+\text{，}\\
\lambda X_\lambda & \text{若 }\lambda\in\RR_+\text{。}
\end{cases}
\]

于是有两种可能：
\begin{itemize}
\item 或者 $B$ 没有任何正实\trEigenvalue{}，此时 $B(1)=-\mathrm{Id}$；
\item 或者 $B$ 有两个正实\trEigenvalue{}，$\lambda$ 与 $1/\lambda$。此时，在诸特征向量 $X_\lambda$ 的（取适当顺序的）\trSymplectic{}基下，矩阵 $B(1)$ 为
\[
B(1)=\left(\begin{array}{cc|c}
\lambda & 0 & \\
0 & 1/\lambda & 0\\
\hline
\multicolumn{2}{c|}{0} & -\mathrm{Id}
\end{array}\right)\text{。}
\]

由于群 $\mathrm{Sp}(2n;\RR)$ 是道路\trConnected{}的，我们可以找到一条道路，把这个矩阵连到\trCanonical{}基下同一形式的矩阵，随后再连到矩阵
\[
\left(\begin{array}{cc|c}
2 & 0 & \\
0 & 1/2 & 0\\
\hline
\multicolumn{2}{c|}{0} & -\mathrm{Id}
\end{array}\right)\text{。}
\]
\end{itemize}
于是，这两个\trOpen{}集 $\mathrm{Sp}(2n)^{\pm}$ 的确就是 $\mathrm{Sp}(2n)^\star$ 的诸\trConnectedComponent。

现在让我们来证明，它们各自到 $\mathrm{Sp}(2n;\RR)$ 中的包含都在\trFundamentalGroup{}上诱导平凡的\trMorphism。正如下面这条引理所容许的那样，只需\trLift{}映射 $\rho$ 即可。

\begin{lemma}\label{thm:rho-lift-diagrams}
存在两个连续映射 $\widetilde{\rho}_\pm\colon\mathrm{Sp}(2n)^{\pm}\to\RR$，使得图表
\[
\begin{tikzcd}[column sep=1.5cm,row sep=1.8em]
&& \RR \ar[d,"\exp"'] \\
\mathrm{Sp}(2n)^{\pm} \ar[r,"i"] \ar[urr,"\widetilde{\rho}_\pm"] & \mathrm{Sp}(2n) \ar[r,"\rho"'] & S^1
\end{tikzcd}
\]
交换。
\end{lemma}

的确，我们已经证明 $\rho$ 在\trFundamentalGroup{}上诱导一个\trIsomorphism。于是，由于
\[
\rho_\star\circ i_\star=\exp_\star\circ(\widetilde{\rho}_\pm)_\star=0\text{，}
\]
\trHomomorphism
\[
i_\star\colon\pi_1\,\mathrm{Sp}(2n)^{\pm}\longrightarrow\pi_1(\mathrm{Sp}(2n))
\]
是零。
\end{proof}

\cref{thm:rho-lift-diagrams}依赖于 $\rho$ 之构造的诸细节。\cref{thm:rho-lift-diagrams}与\cref{thm:distinct-positive-eigenvalues}这两条引理的证明是初等的——初等就初等在它们只用到初等的线性代数。尽管如此，我们觉得把细节写出来还是有用的。它们可以在本章的附录中找到，按理它们就应当在那里，紧跟在 $\rho$ 的构造之后（§7.3.d）。

我们已经完成了第二步；现在只须定义\trIndex{}本身了——一个整数。
